\documentclass[11pt]{article}
\usepackage[margin=1in]{geometry}
\usepackage{amsmath,amssymb,amsthm,mathtools}
\usepackage{empheq}
\usepackage{booktabs}
\usepackage{enumitem}
\usepackage{graphicx}
\usepackage{xcolor}
\usepackage[T1]{fontenc}
\usepackage{titlesec}
\usepackage{authblk}
\usepackage[numbers,sort&compress]{natbib}
\usepackage{hyperref}
\hypersetup{colorlinks=true, linkcolor=blue!45!black, citecolor=blue!45!black,
            urlcolor=blue!45!black}

% ---- SIAM-style run-in section headings ---- %
\titleformat{\section}[runin]{\normalfont\bfseries}{\thesection.}{0.5em}{}[.\quad]
\titleformat{\subsection}[runin]{\normalfont\bfseries}{\thesubsection.}{0.5em}{}[.\quad]
\titleformat{\subsubsection}[runin]{\normalfont\itshape}{\thesubsubsection.}{0.5em}{}[.\quad]
\titlespacing*{\section}{0pt}{1.4em}{0.4em}
\titlespacing*{\subsection}{0pt}{1.0em}{0.4em}
\titlespacing*{\subsubsection}{0pt}{0.7em}{0.4em}

\theoremstyle{plain}
\newtheorem{proposition}{Proposition}
\newtheorem{corollary}{Corollary}
\theoremstyle{definition}
\newtheorem{remark}{Remark}

\newcommand{\TODO}[1]{\textcolor{red}{\textbf{[TODO: #1]}}}
\newcommand{\figph}[2]{%
  \begin{center}\setlength{\fboxsep}{10pt}%
  \fcolorbox{gray!60}{gray!8}{\parbox[c][#1][c]{0.82\linewidth}{%
  \centering\itshape\color{gray!50!black} #2}}\end{center}}

\newcommand{\Lop}{\mathcal{L}}
\newcommand{\Jop}{\mathcal{J}}
\newcommand{\Rop}{\mathcal{R}}
\newcommand{\Ndot}{\dot{N}}
\newcommand{\Nddot}{\ddot{N}}
\newcommand{\Ndddot}{\dddot{N}}
\newcommand{\bO}{\mathcal{O}}
\newcommand{\DT}{\Delta T}
\newcommand{\hf}{h_{\mathrm{f}}}

\title{\bfseries Learning the Truncation Error:\\ A Step-Size Resolution Closure
for Spectral AB2--CN2\\ Integration of Two-Dimensional Quasi-Geostrophic Flows}

\author[1]{Sanaa Mouzahir}
\author[1]{\TODO{co-authors / advisor}}
\affil[1]{\small Department of Mechanical Engineering and Center for
Computational Science and Engineering,\\ Massachusetts Institute of Technology,
Cambridge, MA 02139}
\date{\today}

\begin{document}
\maketitle

\begin{abstract}
\noindent
We develop a learned closure that removes the \emph{time-discretization} error
of the spectral Adams--Bashforth-2 / Crank--Nicolson (AB2--CN2) IMEX
integrator applied to the two-dimensional quasi-geostrophic (QG) vorticity
equation. In contrast to sub-grid (spatial) closures, the scheme and its order
are left unchanged; instead, each coarse step of size $\DT$ is corrected so
that it reproduces the accuracy of the same scheme run at a fine step
$\hf=\DT/K$, or, in the limit, of the exact flow. A modified-equation analysis
gives the local truncation error in closed form to $\bO(\DT^6)$,
$\tau=-\sum_{p=3}^{6}(\DT^p/D_p)\,R_p+\bO(\DT^7)$ with
$D_p=\{12,24,240,1440\}$; all coefficients are verified symbolically and by a
vs-exact-flow convergence study. Each operator $R_p$ splits into an
\emph{analytical} part (powers of the linear operator on $\omega$ and $N$,
diagonal in Fourier space and free at inference) and an \emph{expensive} part
built from time derivatives of the advective Jacobian. A compact,
numerics-informed convolutional network---bilinear (gated) to represent the
Jacobian, with circular padding and a seven-point receptive field---learns only
the expensive part, as a $\DT$-independent target. An empirical magnitude audit
shows that the higher-order operators are governed by the norms of higher time
derivatives of $N$, which grow $\sim\!50\times$ per order in a turbulent
cascade; this gives the modified-equation series a finite radius of convergence
$\DT_\star\!\sim\!2\times10^{-2}$ that bounds the closure's usable step size
independently of correction order. We evaluate on decaying
($\mathrm{Re}\!\sim\!10^5$, $\beta=0$) and forced $\beta$-plane
($\mathrm{Re}\!\sim\!10^4$, $\beta=1$) turbulence.
\TODO{headline result sentence once rollouts are final.}

\smallskip
\noindent\textbf{Key words.} time-integration error, modified equation,
closure modeling, IMEX schemes, pseudo-spectral methods, quasi-geostrophic
turbulence, scientific machine learning
\end{abstract}

% ===================================================================== %
\section{Introduction}
\label{sec:intro}
Numerical simulations of geophysical turbulence are routinely limited by the
time step. In stiff pseudo-spectral solvers the cost is dominated by the
nonlinear (advective) term, evaluated once or twice per step, and the total
cost scales as $T/\Delta t$ to integrate to a horizon $T$. For an integrator
whose accuracy---not linear stability---sets $\Delta t$, halving the usable
step doubles the cost. Reducing this time-discretization cost without changing
the discretization itself is the problem we address.

Data-driven closures are by now standard for representing unresolved
\emph{spatial} scales on a coarse grid, in both engineering large-eddy
simulation and geophysical modeling \citep{gupta2021,guan2022,srinivasan2024}.
The analogous problem for unresolved \emph{temporal} scales---the error
incurred by taking a large time step---has received far less attention. Here we
treat the time-discretization error of a fixed scheme as a closure problem: we
derive the scheme's local truncation error (LTE) in closed form, identify the
single piece that is expensive to evaluate, and learn only that piece.

The construction is a \emph{step-size resolution} closure. The scheme and its
formal order are unchanged; a learned correction makes one coarse step of size
$\DT$ behave as if the same scheme had been run at a finer step $\hf=\DT/K$,
recovering $\hf$-accuracy at $\DT$-cost. Because a fine-step Runge--Kutta
trajectory is the exact flow to machine precision, the $K\!\to\!\infty$ limit
targets the exact flow, and the closure then removes the scheme's full
time-truncation error.

This places the method between three classical ideas while remaining distinct
from each. Modified-equation analysis \citep{warming1974} supplies the
truncation machinery but is normally used to characterize or add order, not to
emulate a finer-step version of the \emph{same} scheme. Defect correction
\citep{stetter1978} iteratively matches a higher-order scheme at the same step;
we instead match the same scheme at a smaller step, a target with no
defect-correction analogue (the gain is set by $K$ at training time, not by a
reference order). Parareal \citep{lions2001} attains fine-propagator accuracy
at coarse cost by iterating in time; our closure is an \emph{amortized}
parareal in which the fine propagator is paid for once during training and
reused at inference without iteration. High-order time integration that pays
once for the curvature of a constrained problem has been developed for
dynamical low-rank systems in the same group \citep{charous2023}.

\paragraph{Notation.}
Table~\ref{tab:notation} collects the notation used throughout.
\begin{table}[h]
\centering\small
\caption{Key notation.}
\label{tab:notation}
\begin{tabular}{ll}
\toprule
symbol & meaning \\
\midrule
$\omega,\ \psi$        & vorticity and streamfunction, $\psi=\Delta^{-1}\omega$ \\
$\Lop$                 & stiff linear operator $\nu\Delta-\mu-\beta\partial_x\Delta^{-1}$ (Fourier-diagonal) \\
$N(\omega)$            & nonlinear term $-\Jop(\psi,\omega)+F$, $F$ time-independent \\
$\Jop(\psi,\omega)$    & advective Jacobian $\psi_x\omega_y-\psi_y\omega_x$ \\
$Q^{(k)},\ \dot Q$     & $k$-th total time derivative of $Q$ along the trajectory \\
$\DT,\ \hf,\ K$        & coarse step, fine step, ratio $K=\DT/\hf$ \\
$\tau$                 & per-step local truncation error \\
$R_p$                  & order-$p$ truncation operator (Prop.~\ref{prop:lte}) \\
$\Rop$                 & leading closure operator $R_3$ \\
$f_{\rm NN}$           & learned (expensive) part of the closure \\
$D_p$                  & denominators $\{12,24,240,1440\}$ for $p=3,\dots,6$ \\
\bottomrule
\end{tabular}
\end{table}

\paragraph{Outline.}
Section~\ref{sec:lte} introduces the spectral AB2--CN2 integrator and derives
its local truncation error to $\bO(\DT^6)$, with a vs-exact-flow convergence
study confirming every operator. Section~\ref{sec:closure} defines the
truncation-error closure, splits each operator into analytical and learned
parts, analyses the magnitude of each part (and the resulting radius of
convergence), and reports the measured magnitudes. Section~\ref{sec:arch}
designs the network that learns the expensive part and documents the
architecture progression that led to it. Section~\ref{sec:scheme} states the
closed AB2--CN2 scheme and its cost. Section~\ref{sec:experiments} reports
numerical experiments on decaying and forced turbulence, and
Section~\ref{sec:conclusion} concludes.

% ===================================================================== %
\section{Local truncation error of spectral AB2--CN2}
\label{sec:lte}

\subsection{The spectral AB2--CN2 integrator}
\label{sec:integrator}
We consider the 2D QG vorticity equation on a doubly-periodic domain,
\begin{equation}
\partial_t\omega = \Lop\,\omega + N(\omega),\qquad
N(\omega)=-\Jop(\psi,\omega)+F,\quad \psi=\Delta^{-1}\omega,
\label{eq:qg}
\end{equation}
with $\Lop=\nu\Delta-\mu-\beta\partial_x\Delta^{-1}$ collecting viscosity,
linear drag, and the $\beta$-plane term. Crank--Nicolson treats the stiff
linear term implicitly and Adams--Bashforth-2 treats the nonlinearity
explicitly:
\begin{equation}
\omega^{n+1}=r(h)\,\omega^n+s(h)\Big(\tfrac32 N^n-\tfrac12 N^{n-1}\Big),\quad
r=\frac{1+\tfrac h2\Lop}{1-\tfrac h2\Lop},\ \ s=\frac{h}{1-\tfrac h2\Lop}.
\label{eq:ab2cn2}
\end{equation}
All linear operators are diagonal in Fourier space, so the implicit solve is an
elementwise division; the only nonlinearity is the Jacobian, evaluated
pseudo-spectrally at five FFTs per call, with $2/3$ dealiasing of the products.
The scheme is second-order accurate in time. Sweeping $\Delta t$ against a fine
reference, started from a common developed-flow snapshot, recovers the expected
$\Delta t^2$ slope (Sec.~\ref{sec:integrator-illus}).

\subsection{Derivation}
\label{sec:lte-derivation}
Along the true trajectory, repeated differentiation of \eqref{eq:qg} gives
$\omega^{(k)}=\Lop^k\omega+\sum_{j=0}^{k-1}\Lop^{k-1-j}N^{(j)}$, and, since
$N=-\Jop(\psi,\omega)+F$ with $F$ static and $\psi^{(k)}=\Delta^{-1}
\omega^{(k)}$, bilinearity of $\Jop$ and the Leibniz rule give the
time-derivatives of $N$,
\begin{equation}
N^{(m)}=-\sum_{j=0}^m\binom{m}{j}\Jop\!\big(\psi^{(m-j)},\omega^{(j)}\big),
\label{eq:leibniz}
\end{equation}
e.g.\ $\Ndot=-\Jop(\dot\psi,\omega)-\Jop(\psi,\dot\omega)$. Equations
\eqref{eq:leibniz} and the $\omega^{(k)}$ recursion are mutually recursive and
are evaluated in interleaved order; each additional time-derivative of $N$
costs another layer of nested Jacobians.

Expanding the exact flow $\omega(t_n+h)=\sum_k\frac{h^k}{k!}\omega^{(k)}$
(equivalently, the Duhamel integral in $\varphi$-functions) and the update
\eqref{eq:ab2cn2} (via the geometric expansions of $r,s$ and the AB2 bracket
$\tfrac32N^n-\tfrac12N^{n-1}$), then matching term by term, gives the LTE.

\begin{proposition}[Local truncation error of AB2--CN2]
\label{prop:lte}
The per-step local truncation error of \eqref{eq:ab2cn2} against the exact flow
of \eqref{eq:qg} is
\begin{equation}
\tau=\omega(t_n+h)-\omega^{n+1}_{\rm AB2CN2}
=-\frac{h^3}{12}R_3-\frac{h^4}{24}R_4-\frac{h^5}{240}R_5
 -\frac{h^6}{1440}R_6+\bO(h^7),
\label{eq:tau}
\end{equation}
with
\begin{align}
R_3 &= \Lop^3\omega+\Lop^2N+\Lop\Ndot-5\Nddot,\label{eq:R3}\\
R_4 &= 2\Lop^4\omega+2\Lop^3N+2\Lop^2\Ndot-4\Lop\Nddot+\Ndddot,\\
R_5 &= 13\Lop^5\omega+13\Lop^4N+13\Lop^3\Ndot-17\Lop^2\Nddot+8\Lop\Ndddot-7N^{(4)},\\
R_6 &= 43\Lop^6\omega+43\Lop^5N+43\Lop^4\Ndot-47\Lop^3\Nddot+28\Lop^2\Ndddot
        -17\Lop N^{(4)}+4N^{(5)}.
\end{align}
The orders $h^0,h^1,h^2$ cancel identically (second-order accuracy).
\end{proposition}

\begin{remark}
The $\Lop^p\omega$ coefficients $\{-\tfrac1{12},-\tfrac1{12},-\tfrac{13}{240},
-\tfrac{43}{1440}\}$ are exactly the error coefficients of the $(1,1)$-Pad\'e
approximant $r(h)$ to $e^{h\Lop}$ (the pure Crank--Nicolson truncation), an
internal consistency check. All coefficients in Prop.~\ref{prop:lte} were
verified symbolically.
\end{remark}

\begin{remark}[$K$-fine reference]
If one targets $K$ fine steps of the same scheme at $\hf=\DT/K$ rather than the
exact flow, every term acquires a factor $(1-1/K^p)$; for the leading term this
is $(1-1/K^2)$, and for $K\gtrsim100$ all such factors are $\approx1$. A
fine-$\hf$ RK4 trajectory equals the exact flow to $\sim10^{-23}$ at
$\DT=10^{-3}$, so the exact-flow operators are what one deploys.
\end{remark}

\subsection{Illustration}
\label{sec:integrator-illus}
We verify Prop.~\ref{prop:lte} numerically by a convergence study against the
exact flow. For a snapshot $\omega^n$ we form the exact history
$\omega^{n-1}=\omega(t_n-h)$ and the exact future $\omega(t_n+h)$ by fine RK4,
take one AB2--CN2 step of size $h$, and measure
$\tau_{\rm emp}=\omega(t_n+h)-\omega^{n+1}$ and the residuals
$\|\tau_{\rm emp}-\tau_{\rm pred}^{(P)}\|$ after subtracting the analytical
prediction kept through order $P$. Keeping $R_3,\dots,R_P$ leaves a residual
$\bO(h^{P+1})$; the measured log--log slopes are $4,5,6$ for $P=3,4,5$,
confirming $R_3,R_4,R_5$ (the slope-7 line for $R_6$ lies at the float64
roundoff floor and is confirmed symbolically instead).
\figph{4.2cm}{Figure placeholder: vs-exact-flow convergence. $\|\tau_{\rm
emp}\|$ (slope 3) and the residuals after subtracting $R_3$, $R_3{+}R_4$,
$R_3{+}R_4{+}R_5$ (slopes 4, 5, 6) vs $h$, with reference-slope guides. Also:
the bare-scheme $\Delta t^2$ convergence of Sec.~\ref{sec:integrator}. Source:
\texttt{validate\_ab2cn2\_vs\_truth.py}.}

% ===================================================================== %
\section{The truncation-error closure}
\label{sec:closure}

\subsection{Derivation}
\label{sec:closure-derivation}
Since $\omega^{n+1}_{\rm AB2CN2}+\tau=\omega(t_n+h)$ exactly, the correction to
\emph{add} to each bare step is $\delta=\tau$. Removing the full truncation
error through $\bO(h^6)$ leaves residual $\bO(h^7)$:
\begin{equation}
\omega^{n+1}_{\rm closed}=\omega^{n+1}_{\rm AB2CN2}
-\frac{\DT^3}{12}R_3-\frac{\DT^4}{24}R_4-\frac{\DT^5}{240}R_5
-\frac{\DT^6}{1440}R_6 .
\label{eq:closed}
\end{equation}
Each operator separates into two qualitatively different parts:
\begin{equation}
R_p=\underbrace{R_p^{\rm anal}}_{\text{powers of }\Lop\text{ on }\omega,N}
   +\underbrace{R_p^{\rm NN}}_{\text{time-derivatives of }N},
\end{equation}
with, for the leading term, $R_3^{\rm anal}=\Lop^3\omega+\Lop^2N$ and
$R_3^{\rm NN}=\Lop\Ndot-5\Nddot$ (Table~\ref{tab:split}). The analytical parts
are diagonal in Fourier space (or use the single Jacobian already needed for
$N$) and are free at inference. The learned parts contain nested time
derivatives of the Jacobian and are the only expensive pieces.
\begin{table}[h]
\centering\small
\caption{Analytical vs.\ learned parts of each operator.}
\label{tab:split}
\begin{tabular}{cll}
\toprule
$p$ & analytical $R_p^{\rm anal}$ & learned $R_p^{\rm NN}$ \\
\midrule
3 & $\Lop^3\omega+\Lop^2N$       & $\Lop\Ndot-5\Nddot$\\
4 & $2\Lop^4\omega+2\Lop^3N$     & $2\Lop^2\Ndot-4\Lop\Nddot+\Ndddot$\\
5 & $13\Lop^5\omega+13\Lop^4N$   & $13\Lop^3\Ndot-17\Lop^2\Nddot+8\Lop\Ndddot-7N^{(4)}$\\
6 & $43\Lop^6\omega+43\Lop^5N$   & $43\Lop^4\Ndot-47\Lop^3\Nddot+28\Lop^2\Ndddot-17\Lop N^{(4)}+4N^{(5)}$\\
\bottomrule
\end{tabular}
\end{table}

\paragraph{Learning target.}
A network learns the leading expensive part as the $\DT$-independent quantity
\begin{equation}
f_{\rm NN}(\omega)=\tfrac{1}{12}\big[\Lop\Ndot-5\Nddot\big],
\label{eq:target}
\end{equation}
which is $\bO(1)$ for QG (no $\DT^3$ cancellation, no tiny labels) and is
reusable across $(\DT,K)$ by analytical rescaling. The deployed correction at
leading order is then
$\delta=-\DT^3(1-1/K^2)\,[\,\tfrac1{12}(\Lop^3\omega+\Lop^2N)+f_{\rm NN}\,]$.

\subsection{Analysis: which orders matter}
\label{sec:analysis}
Whether the higher-order operators contribute depends on the norms
$\|\Ndot\|,\dots,\|N^{(5)}\|$, which are \emph{not} known a priori: each is a
nested Jacobian of an advection field, and in a turbulent cascade the higher
derivatives probe ever finer scales. The per-order LTE contribution is
$\frac{h^p}{D_p}\|R_p\|$. Writing the ratio of consecutive contributions as
$\rho_p(h)=\frac{h^{p+1}/D_{p+1}\,\|R_{p+1}\|}{h^{p}/D_{p}\,\|R_{p}\|}$, the
series converges only while $\rho_p<1$. Because $\rho_p\propto h$, there is a
finite radius of convergence $\DT_\star$ above which consecutive terms stop
shrinking and the modified-equation expansion ceases to represent the error.

\begin{corollary}[Bounded usable step]
\label{cor:radius}
If $\|N^{(p+1)}\|/\|N^{(p)}\|\to g$ (the per-order growth of the nonlinear time
derivatives), then $\DT_\star\sim D_{p+1}/(D_p\,g)\sim 1/g$. The closure's
usable step size is bounded by $\DT_\star$ \emph{regardless of how many
correction orders are carried}; past $\DT_\star$ no finite-order closure
improves accuracy.
\end{corollary}
\noindent The complementary structural bound is the explicit-stability ceiling
of AB2 on the nonlinear term, which limits $\DT$ independently.

\subsection{Illustration}
\label{sec:closure-illus}
We measure the operators directly on developed snapshots, validating the
chain-rule derivatives against RK4 finite differences (to $\le0.7\%$, and
$\le0.6\%$ even for $N^{(5)}$). For forced turbulence
(Table~\ref{tab:mono}), within each $R_p$ the highest time-derivative of $N$
dominates by orders of magnitude; the analytical pieces are negligible. The
pure ($\Lop^0$) norms grow $\sim\!30$--$75\times$ per order
($\|\Nddot\|\!\sim\!7.6\!\times\!10^2$ to
$\|N^{(5)}\|\!\sim\!1.7\!\times\!10^8$), giving $g\!\sim\!50$ and, via
Cor.~\ref{cor:radius}, $\DT_\star\!\sim\!2\!\times\!10^{-2}$
(Table~\ref{tab:lte}).
\begin{table}[h]
\centering\small
\caption{Measured RMS of representative monomials and the assembled
$\|R_p\|$ (mean over snapshots, forced turbulence).}
\label{tab:mono}
\begin{tabular}{lr lr}
\toprule
$R_3$ & RMS & $R_5$ & RMS \\
\midrule
$\Lop^3\omega$ & $1.17\!\times\!10^{-1}$ & $\Lop\Ndddot$ & $3.74\!\times\!10^{4}$\\
$\Lop^2 N$     & $4.59\!\times\!10^{-1}$ & $N^{(4)}$     & $2.27\!\times\!10^{6}$\\
$\Lop\Ndot$    & $1.50\!\times\!10^{1}$  & $\|R_5\|$     & $1.59\!\times\!10^{7}$\\
$\Nddot$       & $7.57\!\times\!10^{2}$  &              & \\
$\|R_3\|$      & $3.78\!\times\!10^{3}$  & $R_6$ & RMS\\
\cmidrule(l){1-2}\cmidrule(l){3-4}
$R_4$ & RMS & $\Lop N^{(4)}$ & $2.59\!\times\!10^{6}$\\
\cmidrule(l){1-2}
$\Lop\Nddot$ & $6.58\!\times\!10^{2}$ & $N^{(5)}$ & $1.70\!\times\!10^{8}$\\
$\Ndddot$    & $3.64\!\times\!10^{4}$ & $\|R_6\|$ & $6.79\!\times\!10^{8}$\\
$\|R_4\|$    & $3.64\!\times\!10^{4}$ &           & \\
\bottomrule
\end{tabular}
\end{table}
\begin{table}[h]
\centering\small
\caption{Measured per-order LTE contribution $\frac{h^p}{D_p}\|R_p\|$ (forced
turbulence). At $\DT=10^{-1}$ the $h^5$ term exceeds the $h^3$ term: the
expansion has diverged.}
\label{tab:lte}
\begin{tabular}{lcccc}
\toprule
$\DT$ & $\tfrac{h^3}{12}\|R_3\|$ & $\tfrac{h^4}{24}\|R_4\|$
      & $\tfrac{h^5}{240}\|R_5\|$ & $\tfrac{h^6}{1440}\|R_6\|$\\
\midrule
$10^{-3}$ & $3.15\!\times\!10^{-7}$ & $1.52\!\times\!10^{-9}$ & $6.63\!\times\!10^{-11}$ & $4.71\!\times\!10^{-13}$\\
$10^{-2}$ & $3.15\!\times\!10^{-4}$ & $1.52\!\times\!10^{-5}$ & $6.63\!\times\!10^{-6}$ & $4.71\!\times\!10^{-7}$\\
$10^{-1}$ & $3.15\!\times\!10^{-1}$ & $1.52\!\times\!10^{-1}$ & $6.63\!\times\!10^{-1}$ & $4.71\!\times\!10^{-1}$\\
\bottomrule
\end{tabular}
\end{table}

\noindent The practical reading: for $\DT\le10^{-3}$, $R_3$ dominates ($R_4$ is
$0.5\%$) and removing $R_3$ alone eliminates the truncation error to the
$\bO(\DT^4)$ floor; near $\DT=10^{-2}$, $R_4$ contributes $\sim5\%$; beyond
$\DT_\star$ the expansion diverges and the closure cannot help. This overturns
the naive expectation that higher-order corrections extend the usable step.
\TODO{add decaying-turbulence column; expect larger $\DT_\star$.}

% ===================================================================== %
\section{Network architecture for the learned operator}
\label{sec:arch}

\subsection{Design}
\label{sec:arch-design}
The target \eqref{eq:target} is dictated by structure, and the architecture is
built to match it.

\subsubsection{Building blocks and locality}
After spatial gradients, $\Jop$ and its $\Lop$-images are \emph{local} (they
are pointwise products of gradient fields, representable by $3\times3$ stencils
and a multiplicative nonlinearity). The only nonlocal operation is
$\Delta^{-1}$, which recovers $\psi$ from $\omega$. To keep the receptive field
small we therefore supply $\psi$ directly rather than asking the network to
invert the Laplacian.

\subsubsection{Minimal input set}
Through the chain rule \eqref{eq:leibniz}, $f_{\rm NN}$ depends on $\omega,\psi$
and their first two time derivatives. Supplying three consecutive time levels
makes those derivatives available by finite differences, giving the minimal,
non-redundant input set
$\big(\omega^n,\omega^{n-1},\omega^{n-2},\psi^n,\psi^{n-1},\psi^{n-2}\big)$
(six channels). The network (\emph{BilinearClosureNet}) then mirrors the
operator: a $1\times1$ stem and time-difference block, a $3\times3$ block with
gated-linear-unit (GLU) activation that represents the bilinear Jacobian, a
residual refinement, a block applying $\Lop$, and a $1\times1$ output
(Table~\ref{tab:arch}). Circular padding throughout matches the periodic
boundary; the receptive field is seven points and the model has
$\approx\!374$k parameters.
\begin{table}[h]
\centering\small
\caption{BilinearClosureNet. $\mathrm{GLU}(x)=a\odot\sigma(b)$ with
$(a,b)=\mathrm{split}(x)$ supplies the bilinear gating that represents $\Jop$;
$\mathrm{GeLU}$ is a smooth nonlinearity.}
\label{tab:arch}
\begin{tabular}{llccc}
\toprule
Block & Operation & Kernel & Activation & Receptive field \\
\midrule
Stem   & channel mix on the 6 inputs           & $1\times1$ & ---  & 1 \\
A      & time finite differences (channel mix) & $1\times1$ & ---  & 1 \\
B      & spatial gradients + bilinear products & $3\times3$ & GLU  & 3 \\
C      & residual refinement                   & $3\times3$ & GeLU & 5 \\
D      & apply $\Lop$ (for the $\Lop\Ndot$ piece)& $3\times3$& ---  & 7 \\
Output & final linear combination              & $1\times1$ & ---  & 7 \\
\bottomrule
\end{tabular}
\end{table}

\subsection{Illustration: architecture progression}
\label{sec:arch-illus}
The minimal design was reached by elimination, each stage isolating one
bottleneck (validation relative-$L^2$ on $f_{\rm NN}$):
\begin{itemize}[leftmargin=1.4em,nosep,topsep=2pt]
\item \textbf{Stage 0 (U-Net, $(\omega_0,\psi_0)$):} $\approx45\%$---no
temporal information, and downsampling washes out the high-$k$ structure where
$\Nddot$ lives.
\item \textbf{Stage 1 (U-Net, gradients $+\omega_{-1}$):} $\approx30\%$.
\item \textbf{Stage 2 (expose analytical $N,\Ndot$):} $\approx28\%$---handing
the network the exact building blocks barely helps, so the bottleneck is
regression, not features.
\item \textbf{Stage 3 (decouple target to $\tfrac1{12}[\Lop\Ndot-5\Nddot]$):}
$\approx25\%$---cleaner target, but the U-Net inductive bias is wrong for a
pointwise bilinear operator.
\item \textbf{Stage 4 (BilinearClosureNet, current best):} add
$\omega_{-2},\psi_{-2}$ and match the architecture to the operator.
\TODO{state final dealiased val.\ rel-$L^2$ ($\approx0.24$) and parameter count.}
\end{itemize}
\figph{3.8cm}{Figure placeholder: architecture-progression loss curves
(Stages 0--4), val.\ rel-$L^2$ vs epoch.}

% ===================================================================== %
\section{The closed AB2--CN2 scheme}
\label{sec:scheme}
At inference, one closed coarse step evaluates the bare update
\eqref{eq:ab2cn2}, the analytical part of $R_3$ (Fourier-diagonal multiplies
plus the cached Jacobian), and one network forward pass, then applies
\eqref{eq:closed} truncated at the order warranted by $\DT$ (Sec.~\ref{sec:analysis}):
\begin{equation}
\omega^{n+1}_{\rm closed}=\omega^{n+1}_{\rm AB2CN2}
-\DT^3\big(1-\tfrac1{K^2}\big)
\Big[\tfrac1{12}\big(\Lop^3\omega^n+\Lop^2N^n\big)+f_{\rm NN}(\omega^n)\Big].
\end{equation}
The network is applied \emph{online}: it is evaluated on the closure
trajectory's own evolving state and its correction feeds the next step, so
errors compound autoregressively and the network never sees truth.
Table~\ref{tab:cost} gives the per-step cost: the closure adds roughly one
batched FFT and one network pass, while requiring $T/\DT$ steps instead of
$T/\hf$---a net $\sim\!K$-fold speedup at $\hf$-accuracy.
\begin{table}[h]
\centering\small
\caption{Per coarse step, in FFT-equivalent units.}
\label{tab:cost}
\begin{tabular}{lcc}
\toprule
 & without closure & with closure \\
\midrule
Jacobians per step            & 1 (+1 cached) & 1 (+1 cached) \\
FFTs for $\Lop^k$ scaffolding & 0 & 0 (spectral multiplies) \\
FFTs for NN inputs            & 0 & $\sim$6 (batchable $\to\!\approx$1) \\
NN forward pass               & 0 & 1 \\
Step for matched accuracy     & $\hf=\DT/K$ & $\DT$ \\
\bottomrule
\end{tabular}
\end{table}

% ===================================================================== %
\section{Numerical experiments}
\label{sec:experiments}
We evaluate on two doubly-periodic configurations spanning the isotropic and
$\beta$-dominated regimes (Table~\ref{tab:cases}). For each, the truth is a
fine ($\hf=10^{-5}$) trajectory; the closure runs at $\DT=10^{-3}$ ($K=100$).
All computations use float64, required because the leading target is
$\bO(\DT^3)\!\approx\!10^{-9}$ before rescaling---below the float32
machine-epsilon floor.
\begin{table}[h]
\centering\small
\caption{Experimental configurations (from the run YAMLs).}
\label{tab:cases}
\begin{tabular}{lll}
\toprule
 & Decaying turbulence & Forced $\beta$-plane turbulence \\
\midrule
Grid                 & $256\times256$       & $512\times512$ \\
Domain               & $2\pi\times2\pi$     & $4\pi\times4\pi$ \\
$\nu$                & $1.025\times10^{-5}$ & $1.025\times10^{-4}$ \\
$\mu$                & $0$                  & $0.02$ \\
$\beta$              & $0$                  & $1.0$ \\
Forcing $F$          & none                 & $-0.1\cos(2x)+0.1\cos(2y)$ \\
$\mathrm{Re}$ (approx.)& $\sim\!10^5$       & $\sim\!10^4$ \\
IC                   & Gaussian, peak $k\in[10,32]$ & Gaussian, peak $k\in[3,5]$ \\
Horizon $T$          & $60$                 & $100$ \\
$\DT$ / $K$          & $10^{-3}$ / $100$    & $10^{-3}$ / $100$ \\
\bottomrule
\end{tabular}
\end{table}

\subsection{Decaying turbulence (\texorpdfstring{$\mathrm{Re}\sim10^5$}{Re~1e5})}
\label{sec:exp-decay}
\TODO{report: single-step closure recovery (predicted vs achieved
rel$_{\rm clos}$/rel$_{\rm bare}$, deviation-cosine); autoregressive rollout
error growth vs the bare coarse run; enstrophy/energy spectra; net speedup.}
\figph{4.6cm}{Figure placeholder: decaying-turbulence rollout. Vorticity fields
(truth / bare / closure) and differences at several times; relative-$L^2$
error growth; enstrophy drift; spectra. Source:
\texttt{rollout\_load\_truth\_compare.py}.}

\subsection{Forced \texorpdfstring{$\beta$}{beta}-plane turbulence (\texorpdfstring{$\mathrm{Re}\sim10^4$}{Re~1e4})}
\label{sec:exp-forced}
\TODO{report the same diagnostics for the forced $\beta$-plane case; relate the
achievable horizon to $\DT_\star$ measured in Sec.~\ref{sec:analysis}.}
\figph{4.6cm}{Figure placeholder: forced-turbulence rollout (same panels).}

\subsection{Achieved recovery versus the true-closure ceiling}
\label{sec:exp-ceiling}
Because the exact correction $R_3$ is available analytically (via the
chain-rule time derivatives of $N$, Sec.~\ref{sec:closure-derivation}), we can separate
two questions that a rollout error alone conflates: \emph{how much error
reduction is attainable} at a given step, and \emph{how much the learned network
actually extracts}. We inject the analytic $R_3$ closure in place of the network
on a fixed $16$-step rollout ladder, over the same six in-distribution
(member, IC) pairs and the same fine-grid truth references used to score the
network. Table~\ref{tab:ceiling} reports the median error reduction over bare
AB2--CN2 for both arms, at three coarse steps.

\begin{table}[h]
\centering\small
\caption{Median error reduction over bare AB2--CN2 on a $16$-step rollout,
in distribution (six (member, IC) pairs; all rungs stable in both arms). The
true-closure column is the attainable ceiling; the network column is the
current learned closure.}
\label{tab:ceiling}
\begin{tabular}{lccc}
\toprule
$\DT$ & Network & True closure ($R_3$) & Fraction of ceiling \\
\midrule
$5\times10^{-3}$   & $3.25\times$ & $59.3\times$ & $5.5\%$ \\
$1\times10^{-2}$   & $8.57\times$ & $29.5\times$ & $29\%$  \\
$1.5\times10^{-2}$ & $4.20\times$ & $89.6\times$ & $4.7\%$ \\
\bottomrule
\end{tabular}
\end{table}

Two features are diagnostic. First, the analytic closure is stable and
strongly corrective at every step, confirming that the $\bO(\DT^4)$-floor
argument (Cor.~\ref{cor:radius}) leaves a large, real signal for
$\DT\le\DT_\star$; the shortfall is entirely in the learned extraction, not in
the theory. Second, the fraction of that ceiling the network recovers is
non-monotone in $\DT$---best near the middle of the sweep and worst at both
ends. This is the signature expected of a \emph{single} step-independent stencil
fitting the pooled mean of a target correction whose amplitude scales as
$(\DT\,\sigma)^{S-m}$ (Sec.~\ref{sec:analysis}): one stencil is a compromise
that is accurate near the pooled operating point and degrades toward the sweep
extremes. A spectral decomposition localizes the residual at low wavenumber
(regime-scale structures). Both observations motivate the conditioned stencil of
Sec.~\ref{sec:conclusion}, in which the analytic $\DT^{S-m}$ prefactor is
factored out and only a regime-dependent gain $g_\theta(k;\mathrm{Re},\beta,\mu)$
is learned. \TODO{replace the ``Network'' column once the conditioned-stencil
run lands; report the closed gap.}

% ===================================================================== %
\section{Conclusion}
\label{sec:conclusion}
We have treated the time-discretization error of the spectral AB2--CN2
integrator as a closure problem and removed it with a learned correction
derived from modified-equation analysis. The local truncation error is known
in closed form to $\bO(\DT^6)$ (Prop.~\ref{prop:lte}), every operator verified
symbolically and by a vs-exact-flow convergence study; each operator splits
into a free analytical part and an expensive part built from time derivatives
of the Jacobian, the latter learned by a compact, structure-matched network.
An empirical magnitude audit shows that the expensive parts are governed by the
rapidly growing norms of higher time derivatives of $N$, which gives the
modified-equation series a finite radius of convergence $\DT_\star$ that
bounds the closure's usable step independently of correction order
(Cor.~\ref{cor:radius})---so in practice $R_3$ is the operative closure and
removing it eliminates the truncation error up to the $\bO(\DT^4)$ floor for
$\DT\le\DT_\star$.

\paragraph{Future directions.}
The headroom of Sec.~\ref{sec:exp-ceiling}---a large, stable, analytically
available signal that the current single stencil extracts only in part, most
severely at the sweep extremes---points directly to a \emph{physics-conditioned}
stencil. The error analysis predicts an optimal correction of the form
$\delta_\theta(k)=\DT^{S-m}\,[\text{analytic}]\times g_\theta(k;\mathrm{Re},
\beta,\mu)$: the step-size dependence is factored out in closed form and only a
regime-dependent spectral gain is learned, removing the pooled-mean compromise
that a step-independent stencil is forced into. This makes $(\DT,\mathrm{Re},
\beta,\mu)$ explicit features and is the natural route to cross-regime transfer.
Further directions: out-of-distribution step/horizon tests; learning additional
$N$-derivative orders only where the magnitude audit warrants; and obstacle
flows. \TODO{fold in the conditioned-stencil results once landed.}

\section*{Reproducibility}
The spectral QG solver follows \citet{sureshbabu2025}. Truncation operators are
verified by \texttt{validate\_ab2cn2\_vs\_truth.py}; magnitudes by
\texttt{measure\_truncation\_magnitudes.py}; rollouts by
\texttt{rollout\_load\_truth\_compare.py}. \TODO{code/data availability + repo.}

\bibliographystyle{plainnat}
\begin{thebibliography}{99}
\bibitem[Suresh Babu et~al.(2025)]{sureshbabu2025}
A.~N. Suresh Babu, A.~Sadam, and P.~F.~J. Lermusiaux. \TODO{title}. 2025.
\bibitem[Charous \& Lermusiaux(2023)]{charous2023}
A.~Charous and P.~F.~J. Lermusiaux. Stable rank-adaptive dynamically orthogonal
Runge--Kutta schemes. \emph{SIAM J.\ Sci.\ Comput.}, 2023.
\bibitem[Gupta \& Lermusiaux(2021)]{gupta2021}
A.~Gupta and P.~F.~J. Lermusiaux. Neural closure models for dynamical systems.
\emph{Proc.\ R.\ Soc.\ A}, 477:20201004, 2021.
\bibitem[Guan et~al.(2022)]{guan2022}
Y.~Guan, A.~Chattopadhyay, A.~Subel, and P.~Hassanzadeh. Stable a posteriori
LES of 2D turbulence using convolutional neural networks. \emph{J.\ Comput.\
Phys.}, 458:111090, 2022.
\bibitem[Srinivasan et~al.(2024)]{srinivasan2024}
K.~Srinivasan, M.~D. Chekroun, and J.~C. McWilliams. \TODO{title}. 2024.
\bibitem[Warming \& Hyett(1974)]{warming1974}
R.~F. Warming and B.~J. Hyett. The modified equation approach to the stability
and accuracy analysis of finite-difference methods. \emph{J.\ Comput.\ Phys.},
14:159--179, 1974.
\bibitem[Stetter(1978)]{stetter1978}
H.~J. Stetter. The defect correction principle and discretization methods.
\emph{Numer.\ Math.}, 29:425--443, 1978.
\bibitem[Lions et~al.(2001)]{lions2001}
J.-L. Lions, Y.~Maday, and G.~Turinici. R\'esolution d'EDP par un sch\'ema en
temps ``parar\'eel''. \emph{C.\ R.\ Acad.\ Sci.\ Paris}, 332:661--668, 2001.
\TODO{add: Kraichnan/Leith 2D turbulence; Crank--Nicolson; Adams--Bashforth.}
\end{thebibliography}

\end{document}
